\documentclass[10pt, a4paper]{article}

\usepackage[T1]{fontenc}
\usepackage[utf8]{inputenc}
\usepackage{tgpagella}
\usepackage{microtype}
\usepackage[spanish, es-noshorthands]{babel}
\usepackage{amsmath, amssymb, bm}
\usepackage{multicol}
\usepackage[most]{tcolorbox}
\usepackage[margin=1.5cm]{geometry}
\usepackage{fancyhdr}

\pagestyle{fancy}
\fancyhf{}
\setlength{\headheight}{16pt}
\lhead{\textbf{IMT-342} | Robótica}
\chead{Formulario Oficial — EC2 Diferido}
\rhead{Unidad II}
\lfoot{Material Autorizado}
\rfoot{Página \thepage}
\renewcommand{\headrulewidth}{0.4pt}
\renewcommand{\footrulewidth}{0.4pt}

\definecolor{ucbblue}{RGB}{0, 51, 102}

\tcbset{
    sectionbox/.style={
        colback=ucbblue!5,
        colframe=ucbblue,
        fonttitle=\bfseries,
        left=2mm, right=2mm, top=1mm, bottom=1mm,
        boxsep=0mm,
        arc=1mm
    }
}

\begin{document}

\begin{center}
    \Large\textbf{\textcolor{ucbblue}{Formulario Oficial — EC2 Diferido}}\\[0.2cm]
\end{center}

\begin{tcolorbox}[colback=yellow!10, colframe=black, title=\textbf{Aviso y Reglas de Uso}[cite: 11], left=2mm, right=2mm, top=1mm, bottom=1mm]
    \small
    Este formulario proporciona relaciones matemáticas de la Unidad II. No indica qué fórmula corresponde a cada apartado ni contiene procedimientos resueltos[cite: 11]. \\
    \textbf{Permitido:} esta hoja + calculadora científica no programable[cite: 11]. \\
    \textbf{No permitido:} apuntes, CAS, laptop, Python/Jupyter, teléfono o material adicional[cite: 11].
\end{tcolorbox}

\vspace{0.2cm}
\begin{multicols*}{2}
\small

% --- Transformaciones homogéneas ---
\begin{tcolorbox}[sectionbox, title=Transformaciones Homogéneas{[cite: 11]}]
\begin{equation*}
    T= \begin{bmatrix} R & p \\ 0 & 1 \end{bmatrix} \qquad T^{-1} = \begin{bmatrix} R^T & -R^T p \\ 0 & 1 \end{bmatrix}
\end{equation*}
$R^{-1}=R^T$ (para una matriz de rotación válida)[cite: 11].
\end{tcolorbox}

% --- Cinemática inversa ---
\begin{tcolorbox}[sectionbox, title=Cinemática Inversa{[cite: 11]}]
\begin{equation*}
    p_W = p_6 - d_6 z_6 \qquad z_6 = R_6 \begin{bmatrix} 0 \\ 0 \\ 1 \end{bmatrix}
\end{equation*}
\begin{equation*}
    R_3^6 = ({}^0R_3)^T R_0^6 \qquad s = \pm\sqrt{1-c^2}
\end{equation*}
\begin{equation*}
    q = \operatorname{atan2}(s,c)
\end{equation*}
Verificación genérica: $\boxed{ T(\mathbf{q}^{(k)}) \approx T_d }$[cite: 11].
\end{tcolorbox}

% --- Geometría vectorial ---
\begin{tcolorbox}[sectionbox, title=Geometría Vectorial{[cite: 11]}]
\begin{equation*}
    a \times b = \begin{bmatrix} a_y b_z - a_z b_y \\ a_z b_x - a_x b_z \\ a_x b_y - a_y b_x \end{bmatrix}
\end{equation*}
Producto punto: $a^T b = a_x b_x + a_y b_y + a_z b_z$[cite: 11]. \\
Vectores perpendiculares: $a^T b = 0$[cite: 11].
\end{tcolorbox}

% --- Jacobiano geométrico ---
\begin{tcolorbox}[sectionbox, title=Jacobiano Geométrico{[cite: 11]}]
\begin{equation*}
    v_e = \begin{bmatrix} v_L \\ \omega \end{bmatrix} = J(\mathbf{q})\dot{\mathbf{q}}
\end{equation*}
Articulación revoluta:
\begin{equation*}
    J_i = \begin{bmatrix} z_{i-1} \times (o_n - o_{i-1}) \\ z_{i-1} \end{bmatrix}
\end{equation*}
Articulación prismática:
\begin{equation*}
    J_i = \begin{bmatrix} z_{i-1} \\ 0 \end{bmatrix}
\end{equation*}
\end{tcolorbox}

% --- Rango y singularidad ---
\begin{tcolorbox}[sectionbox, title=Rango y Singularidad{[cite: 11]}]
Para $J \in \mathbb{R}^{m \times n}$: $\operatorname{rank}_{\max}J = \min(m,n)$[cite: 11]. \\
Criterio general de singularidad:
\begin{equation*}
    \boxed{\operatorname{rank}J < \operatorname{rank}_{\max}}
\end{equation*}
Para matriz cuadrada, $\det J = 0$ implica pérdida de rango completo[cite: 11].
\begin{equation*}
    \det \begin{bmatrix} a & b \\ c & d \end{bmatrix} = ad - bc
\end{equation*}
\end{tcolorbox}

% --- LSPB ---
\begin{tcolorbox}[sectionbox, title=LSPB (Vel. Inicial/Final Nula){[cite: 11]}]
$\Delta q = q_f - q_i$. \\
Aceleración mínima: $\boxed{ |\ddot{q}_c| \ge \frac{4|\Delta q|}{t_f^2} }$[cite: 11]. \\
Tiempo de blend (signos consistentes):
\begin{equation*}
    t_c = \frac{t_f}{2} - \frac{1}{2} \sqrt{ t_f^2 - \frac{4(q_f - q_i)}{\ddot{q}_c} }
\end{equation*}
Velocidad de crucero: $\dot{q}_c = \ddot{q}_c t_c$[cite: 11]. \\
\textbf{Fase Aceleración:}
\begin{equation*}
    q(t) = q_i + \frac{1}{2}\ddot{q}_c t^2 \quad \dot{q}(t) = \ddot{q}_c t \quad \ddot{q}(t) = \ddot{q}_c
\end{equation*}
\textbf{Fase Vel. Constante:}
\begin{equation*}
    q(t) = q_i + \ddot{q}_c t_c \left( t - \frac{t_c}{2} \right) \quad \dot{q}(t) = \dot{q}_c \quad \ddot{q}(t) = 0
\end{equation*}
\textbf{Fase Deceleración:}
\begin{align*}
    q(t) &= q_f - \frac{1}{2}\ddot{q}_c (t_f - t)^2 \\
    \dot{q}(t) &= \ddot{q}_c (t_f - t) \quad \ddot{q}(t) = -\ddot{q}_c
\end{align*}
\end{tcolorbox}

% --- Energía y Dinámica ---
\begin{tcolorbox}[sectionbox, title=Energía y Lagrange{[cite: 11]}]
\begin{equation*}
    L = T - V \qquad \frac{d}{dt} \left( \frac{\partial L}{\partial \dot{q}_i} \right) - \frac{\partial L}{\partial q_i} = \tau_i
\end{equation*}
Cinética Rotacional: $T = \frac{1}{2}J\omega^2$[cite: 11]. \\
Cinética Traslacional: $T = \frac{1}{2}mv^2$[cite: 11]. \\
Potencial Gravitacional: $V = mgh$[cite: 11].
\end{tcolorbox}

\begin{tcolorbox}[sectionbox, title=Dinámica General y Modelo 1R{[cite: 11]}]
Dinámica General de Manipuladores:
\begin{equation*}
    \tau = M(\mathbf{q})\ddot{\mathbf{q}} + C(\mathbf{q},\dot{\mathbf{q}})\dot{\mathbf{q}} + \mathbf{g}(\mathbf{q})
\end{equation*}
Modelo 1R (Clase):
\begin{align*}
    J &= I_c + m\ell_c^2 \quad &T &= \frac{1}{2}J\dot{q}^2 \\
    V &= mg\ell_c(1 - \cos q) \quad &\tau &= J\ddot{q} + mg\ell_c\sin q
\end{align*}
\end{tcolorbox}

% --- Trigonometría e Identidades ---
\begin{tcolorbox}[sectionbox, title=Identidades y Unidades{[cite: 11]}]
\begin{align*}
    &\sin^2\theta + \cos^2\theta = 1 \\
    &\sin(A \pm B) = \sin A\cos B \pm \cos A\sin B \\
    &\cos(A \pm B) = \cos A\cos B \mp \sin A\sin B \\
    &\sin(-\theta) = -\sin\theta \qquad \cos(-\theta) = \cos\theta
\end{align*}
Unidades: $[\dot{q}]$=rad/s, $[\ddot{q}]$=rad/s$^2$, $[\tau]$=N$\cdot$m, $[J_{\text{in}}]$=kg$\cdot$m$^2$[cite: 11].
\end{tcolorbox}

\end{multicols*}
\end{document}
